% GENERATED FILE. Edit canonical lessons, then run npm run generate:books.
% canonical-source-hash: fnv1a64:5ff6c98ecc1341a1
% canonical-lessons: UR-C433-tila, UR-C433-dher, UR-C433-ghar, UR-C433-tinka, UR-C433-ukharna

\chapter{Mounds, Heaps and Caves}
\label{ch:ur-mounds-heaps-and-caves}
\hlchaptermodality{433}

\begin{chapteropening}
\textbf{By the end of this chapter:} \emph{I can describe the ground: dunes, heaps, caves and what a storm leaves behind.}

Complete the last lesson of chapter 433: I can describe the ground: dunes, heaps, caves and what a storm leaves behind.
\end{chapteropening}

\section[\texorpdfstring{\ur{ٹیلا} (ṭīlā)}{ٹیلا (ṭīlā)}]{\ur{ٹیلا} (ṭīlā) — a mound; a dune}

\label{lesson:UR-C433-tila}



\begin{quote}
Before the new one: say the Urdu for the glow of sunset, then the Urdu for daylight.
\end{quote}

\subsection*{You'll want to know: \ur{ٹیلا}}
\textbf{\ur{ٹیلا}} — \emph{ṭīlā} — \textquotedblleft{}a mound; a dune\textquotedblright{}.

\begin{quote}
\ur{ریگستان} \ur{میں} \ur{ریت} \ur{کے} \ur{ٹیلے} \ur{ہیں۔}
\end{quote}

\begin{quote}
\emph{registān meṅ ret ke ṭīle haiṅ.}
\end{quote}

\begin{quote}
There are sand dunes in the desert.
\end{quote}

Masculine.

\begin{cousinweb}
\textbf{\ur{ٹیلا}} (\emph{ṭīlā}) is inherited: CDIAL lists the Hindi cousin \emph{ṭīlā}, \textquotedblleft{}a hillock, a mound\textquotedblright{}, under \emph{ṭalla} and its variant \emph{ṭilla}, \textquotedblleft{}a heap\textquotedblright{} (CDIAL 5451).
\end{cousinweb}

\begin{grammarlens}[title={What you've built}]
One more word for getting around.
\end{grammarlens}

\subsection*{Your turn}
\begin{itemize}
\raggedright
  \item \emph{Say it:} \emph{ṭīlā}
  \item \emph{Say it:} \emph{ṭīlā}, once more
  \item \emph{Say it:} \emph{registān meṅ ret ke ṭīle haiṅ.} — There are sand dunes in the desert.
  \item \emph{Recall:} say the Urdu for the glow of sunset, then the Urdu for daylight, then say \emph{ṭīlā} again
\end{itemize}

\begin{tcolorbox}[breakable,colback=teal!4,colframe=teal!35!black,title={Before you move on}]
What does \ur{ٹیلا} mean? (\textquotedblleft{}A mound; a dune\textquotedblright{}.) Say it once more.
\end{tcolorbox}
\section[\texorpdfstring{\ur{ڈھیر} (ḍher)}{ڈھیر (ḍher)}]{\ur{ڈھیر} (ḍher) — a heap; a pile}

\label{lesson:UR-C433-dher}



\begin{quote}
Before the new one: say the Urdu for daylight, then the Urdu for a mound.
\end{quote}

\subsection*{You'll want to know: \ur{ڈھیر}}
\textbf{\ur{ڈھیر}} — \emph{ḍher} — \textquotedblleft{}a heap; a pile\textquotedblright{}.

\begin{quote}
\ur{کسان} \ur{کے} \ur{پاس} \ur{گندم} \ur{کا} \ur{ڈھیر} \ur{ہے۔}
\end{quote}

\begin{quote}
\emph{kisān ke pās gandum kā ḍher hai.}
\end{quote}

\begin{quote}
The farmer has a heap of wheat.
\end{quote}

Masculine.

\begin{cousinweb}
\textbf{\ur{ڈھیر}} (\emph{ḍher}) is inherited: CDIAL lists the Hindi cousin \emph{ḍher}, \textquotedblleft{}a heap\textquotedblright{}, under \emph{ḍhēra}, a form it reconstructs (CDIAL 5599).
\end{cousinweb}

\begin{grammarlens}[title={What you've built}]
One more word for getting around.
\end{grammarlens}

\subsection*{Your turn}
\begin{itemize}
\raggedright
  \item \emph{Say it:} \emph{ḍher}
  \item \emph{Say it:} \emph{ḍher}, once more
  \item \emph{Say it:} \emph{kisān ke pās gandum kā ḍher hai.} — The farmer has a heap of wheat.
  \item \emph{Recall:} say the Urdu for daylight, then the Urdu for a mound, then say \emph{ḍher} again
\end{itemize}

\begin{tcolorbox}[breakable,colback=teal!4,colframe=teal!35!black,title={Before you move on}]
What does \ur{ڈھیر} mean? (\textquotedblleft{}A heap; a pile\textquotedblright{}.) Say it once more.
\end{tcolorbox}
\section[\texorpdfstring{\ur{غار} (ghār)}{غار (ghār)}]{\ur{غار} (ghār) — a cave}

\label{lesson:UR-C433-ghar}



\begin{quote}
Before the new one: say the Urdu for a mound, then the Urdu for a heap.
\end{quote}

\subsection*{You'll want to know: \ur{غار}}
\textbf{\ur{غار}} — \emph{ghār} — \textquotedblleft{}a cave\textquotedblright{}.

\begin{quote}
\ur{پہاڑ} \ur{میں} \ur{ایک} \ur{پرانا} \ur{غار} \ur{ہے۔}
\end{quote}

\begin{quote}
\emph{pahāṛ meṅ ek purānā ghār hai.}
\end{quote}

\begin{quote}
There is an old cave in the mountain.
\end{quote}

Masculine.

\begin{cousinweb}
\textbf{\ur{غار}} (\emph{ghār}) is usually traced, through Persian, to Arabic \emph{ghār}, \textquotedblleft{}a cave, a grotto\textquotedblright{}.
\end{cousinweb}

\begin{grammarlens}[title={What you've built}]
One more word for getting around.
\end{grammarlens}

\subsection*{Your turn}
\begin{itemize}
\raggedright
  \item \emph{Say it:} \emph{ghār}
  \item \emph{Say it:} \emph{ghār}, once more
  \item \emph{Say it:} \emph{pahāṛ meṅ ek purānā ghār hai.} — There is an old cave in the mountain.
  \item \emph{Recall:} say the Urdu for a mound, then the Urdu for a heap, then say \emph{ghār} again
\end{itemize}

\begin{tcolorbox}[breakable,colback=teal!4,colframe=teal!35!black,title={Before you move on}]
What does \ur{غار} mean? (\textquotedblleft{}A cave\textquotedblright{}.) Say it once more.
\end{tcolorbox}
\section[\texorpdfstring{\ur{تنکا} (tinkā)}{تنکا (tinkā)}]{\ur{تنکا} (tinkā) — a straw; a blade of grass}

\label{lesson:UR-C433-tinka}



\begin{quote}
Before the new one: say the Urdu for a heap, then the Urdu for a cave.
\end{quote}

\subsection*{You'll want to know: \ur{تنکا}}
\textbf{\ur{تنکا}} — \emph{tinkā} — \textquotedblleft{}a straw; a blade of grass\textquotedblright{}.

\begin{quote}
\ur{ہوا} \ur{میں} \ur{تنکے} \ur{اڑ} \ur{رہے} \ur{تھے۔}
\end{quote}

\begin{quote}
\emph{hawā meṅ tinke uṛ rahe the.}
\end{quote}

\begin{quote}
Bits of straw were blowing about in the wind.
\end{quote}

Masculine.

\begin{cousinweb}
\textbf{\ur{تنکا}} (\emph{tinkā}) is inherited: CDIAL lists the Hindi cousin \emph{tinkā}, \textquotedblleft{}a blade of grass, a straw\textquotedblright{}, under Sanskrit \emph{tṛṇa}, \textquotedblleft{}grass\textquotedblright{} (CDIAL 5906).
\end{cousinweb}

\begin{grammarlens}[title={What you've built}]
One more word for getting around.
\end{grammarlens}

\subsection*{Your turn}
\begin{itemize}
\raggedright
  \item \emph{Say it:} \emph{tinkā}
  \item \emph{Say it:} \emph{tinkā}, once more
  \item \emph{Say it:} \emph{hawā meṅ tinke uṛ rahe the.} — Bits of straw were blowing about in the wind.
  \item \emph{Recall:} say the Urdu for a heap, then the Urdu for a cave, then say \emph{tinkā} again
\end{itemize}

\begin{tcolorbox}[breakable,colback=teal!4,colframe=teal!35!black,title={Before you move on}]
What does \ur{تنکا} mean? (\textquotedblleft{}A straw; a blade of grass\textquotedblright{}.) Say it once more.
\end{tcolorbox}
\section[\texorpdfstring{\ur{اکھڑنا} (ukhaṛnā)}{اکھڑنا (ukhaṛnā)}]{\ur{اکھڑنا} (ukhaṛnā) — to be uprooted; to come loose}

\label{lesson:UR-C433-ukharna}



\begin{quote}
Before the new one: say the Urdu for a cave, then the Urdu for a straw.
\end{quote}

\subsection*{You'll want to know: \ur{اکھڑنا}}
\textbf{\ur{اکھڑنا}} — \emph{ukhaṛnā} — \textquotedblleft{}to be uprooted; to come loose\textquotedblright{}.

\begin{quote}
\ur{طوفان} \ur{میں} \ur{ایک} \ur{پرانا} \ur{درخت} \ur{اکھڑ} \ur{گیا۔}
\end{quote}

\begin{quote}
\emph{tūfān meṅ ek purānā darakht ukhaṛ gayā.}
\end{quote}

\begin{quote}
An old tree was uprooted in the storm.
\end{quote}

Past (it …): \ur{وہ} \ur{پودا} \ur{اکھڑ} \ur{گیا}, \emph{voh paudā ukhaṛ gayā}.

\begin{cousinweb}
\textbf{\ur{اکھڑنا}} (\emph{ukhaṛnā}) is inherited: CDIAL lists the Hindi cousin \emph{ukhaṛnā}, \textquotedblleft{}to be rooted up\textquotedblright{}, under \emph{utskarati}, a form it reconstructs (CDIAL 1895).
\end{cousinweb}

\begin{grammarlens}[title={What you've built}]
One more word for getting around.
\end{grammarlens}

\subsection*{Your turn}
\begin{itemize}
\raggedright
  \item \emph{Say it:} \emph{ukhaṛnā}
  \item \emph{Say it:} \emph{ukhaṛnā}, once more
  \item \emph{Say it:} \emph{tūfān meṅ ek purānā darakht ukhaṛ gayā.} — An old tree was uprooted in the storm.
  \item \emph{Recall:} say the Urdu for a cave, then the Urdu for a straw, then say \emph{ukhaṛnā} again
\end{itemize}

\begin{tcolorbox}[breakable,colback=teal!4,colframe=teal!35!black,title={Before you move on}]
What does \ur{اکھڑنا} mean? (\textquotedblleft{}To be uprooted; to come loose\textquotedblright{}.) Say it once more.
\end{tcolorbox}
